% 第1節 集合と命題
\section{集合と命題}

\subsection{集合}

\begin{definitionbox}[def:集合]{\textbf{集合}}
ものの集まりを集合という.
\end{definitionbox}

\begin{theorembox}[thm:ド・モルガン]{\textbf{ド・モルガンの法則}}
全体集合を \(U\) とし, \(A,B \subset U\) とする. このとき
\begin{align}
  (A \cup B)^{c} &= \answermath{A^{c} \cap B^{c}} \\
  (A \cap B)^{c} &= \answermath{A^{c} \cup B^{c}}
\end{align}
が成り立つ.
\end{theorembox}

\newpage

\subsection{命題}

\begin{definitionbox}[def:命題]{\textbf{命題}}
真または偽のいずれかが定まる文を命題という.
\end{definitionbox}

